\documentclass[11pt,a4paper]{article}
\usepackage[T1]{fontenc}
\usepackage[utf8]{inputenc}
\usepackage{amsmath,amssymb,xcolor}
\usepackage[margin=25mm]{geometry}
\usepackage[colorlinks=true,urlcolor=blue]{hyperref}
\definecolor{nexorablue}{HTML}{2459B5}
\setlength{\parindent}{0pt}
\setlength{\parskip}{7pt}
\setlength{\emergencystretch}{2em}
\renewcommand{\familydefault}{\sfdefault}
\begin{document}

{\small\bfseries\color{nexorablue} NEXORA CPGE}\par
{\small MPSI / NEX-CIR-01 / Énoncé}\par
{\LARGE\bfseries Propagation et courants signés}\par
\vspace{0.25em}\hrule\vspace{0.6em}

Sur une carte de taille $15\,\mathrm{cm}$, prendre $v=2{,}0\times10^8\,\mathrm{m\,s^{-1}}$. Comparer le temps de propagation aux périodes de signaux de $2$ et $500\,\mathrm{MHz}$. À basse fréquence, deux flèches entrantes d’un nœud portent $3{,}5$ et $-1{,}2\,\mathrm{mA}$. Une flèche sortante porte $0{,}8\,\mathrm{mA}$. Déterminer le courant orienté suivant la quatrième flèche, sortante.


\end{document}
